\begin{textsample}{POS Dim 8 – global\_south\_2025\_01 – Score 8.00 – global\_south\_2025\_01/120626790.txt}  \label{ex:f8_pos_115}
When Israeli troops withdrew from Rafah -- marking the start of the ceasefire agreement earlier this month -- they left behind a trail of destruction in Gaza ' s southernmost city.
An estimated 1.3 million people had sought refuge in Rafah after being displaced from other areas of Gaza during the genocidal war.
All have gone through harrowing experiences.
Ahmed al-Sufi, the city ' s mayor, stated during a recent press conference that the authorities did not have the capacity to deal with the devastation caused by Israel.
As well as subjecting Rafah to numerous airstrikes over the past 15 months, Israeli troops undertook a ground invasion of the city in May."The infrastructure is completely paralyzed,"al-Sufi said."Water and electricity networks are down and the main roads are unusable.
Thousands of homes and public facilities no longer exist."Before Israel declared its genocidal war in October 2023, Rafah had a population of approximately 275,000.
Fatima Abu Taha lived in the Saudi area, west of the @ @ @ @ @ @ @ @ @ @ Israel ' s brutality."I couldn't identify my house until my neighbors helped me,"she told The \textbf{Electronic} Intifada."The entire neighborhood has been turned into wasteland, without any buildings.
Even the desalination plant we relied on for water has been completely destroyed."Although the Israeli military has withdrawn from the center of Rafah, its troops have maintained a presence just a few hundred meters from Fatima Abu Taha ' s destroyed home."Even staying in a tent on the ruins of my old home feels like a great risk,"she said.
Despite the ceasefire, Israeli forces have opened fire on and killed Palestinians attempting to return to their homes in Rafah."Everything is gone"Adham Mansour was completely shocked when he went back to the area around Barqa Stadium in the northern part of Rafah.
Before his \textbf{displacement}, Mansour lived and owned a clothing store in the area."I came to check on our five-story home, but I @ @ @ @ @ @ @ @ @ @"It ' s as if nothing ever stood here.
Even my clothing store, my only source of income, has been completely leveled.""Everything is gone,"he said."They have destroyed my entire life.
I don't even know where to start after losing my home and livelihood all at once."Ahmad Farhat returned to what remained of his home in the Tal al-Sultan area of Rafah."I consider myself lucky because my house wasn't completely destroyed,"he said."But it was hit by shells that caused significant damage.""My home is uninhabitable now,"he added."I will need to repair it but the real challenge is getting building materials.
Israel was already restricting their entry before the war.
I can't imagine how long it will take now with such widespread destruction."Ahmad Adwan is an engineer, specializing in damage assessment."The level of devastation here is terrifying,"@ @ @ @ @ @ @ @ @ @ months just to document the damage before reconstruction can even begin."The task of rebuilding Rafah is further complicated by how the amount of available land has shrunk.
As the genocide got underway in October 2023, Eli Cohen, then Israel ' s foreign minister, stated that the territory of Gaza would decrease.
Israel has carried out that threat by expanding the so-called"buffer zone"from which Palestinians are forbidden.
To expand the zone, Israel has laid waste to large tracts of farmland and demolished a huge number of buildings.
The expansion of the zone"makes life here nearly impossible,"said Ahmad Adwan, the engineer.
He added that"the city needs urgent international intervention"for reconstruction as"local resources are nearly nonexistent."Name : Email : The ceasefire announcement earlier this month brought joy to Gaza.
Yet in Rafah, that joy was tempered by shock, particularly once people saw how their homes had been destroyed or badly damaged and @ @ @ @ @ @ @ @ @ @ Posts Bern ( Quds News Network ) - Palestinian-American journalist Ali Abunimah has confirmed that Swiss authorities released and deported him after detaining him for three days over his advocacy for Palestinian rights.
Abunimah,...
Support Countercurrents Countercurrents.org is on annual \textbf{fundraising} drive.
CC is 100 \% \textbf{reader} supported. \textbf{Kindly} \textbf{lend} your helping hand to keep CC alive.
Don't wait for others to do the job for you. \textbf{Please} \textbf{chip} in with whatever you can.

% matched lemmas: chip, displacement, electronic, fundraising, kindly, lend, please, reader
\end{textsample}
